\documentclass[11pt,a4paper]{article}

\usepackage[T1]{fontenc}
\usepackage[utf8]{inputenc}
\usepackage{lmodern}
\usepackage{amssymb,amsmath,amsthm,mathtools}
\usepackage[margin=2.5cm]{geometry}
\usepackage{parskip}
\usepackage[hidelinks]{hyperref}
\usepackage{microtype}

% Theorem environments (arXiv / journal style)
\theoremstyle{plain}
\newtheorem{theorem}{Theorem}[section]
\newtheorem{proposition}[theorem]{Proposition}
\newtheorem{lemma}[theorem]{Lemma}
\newtheorem{corollary}[theorem]{Corollary}

\theoremstyle{definition}
\newtheorem{definition}[theorem]{Definition}
\newtheorem{example}[theorem]{Example}
\newtheorem{remark}[theorem]{Remark}

% Notation
\newcommand{\Shape}{\mathrm{Shape}}
\newcommand{\Pos}{\mathrm{Pos}}
\newcommand{\Type}{\mathbf{Type}}
\newcommand{\id}{\mathrm{id}}
\newcommand{\src}{\mathrm{src}}
\newcommand{\tgt}{\mathrm{tgt}}
\newcommand{\rootpos}{\mathrm{root}}
\newcommand{\sub}{\mathrm{sub}}
\newcommand{\shift}{\mathrm{shift}}
\newcommand{\N}{\mathbb{N}}

\title{The DAG and Graph Rungs:\\
  Failure of Thinness and Finiteness}
\author{Extracted from a machine-checked development}
\date{}

\begin{document}
\maketitle

\begin{abstract}
The network directed container of a polytree is thin and has finite position
sets whenever the underlying graph is finite and acyclic. Both conclusions
fail as soon as one leaves the polytree rung of the structural ladder. We
exhibit two elementary counterexamples. First, extending a wine-shipping
graph by a Singapore transshipment produces a directed acyclic graph with two
distinct routes between the same ports, so the network is not thin. Second, a
two-node scrap-metal cycle yields infinitely many positions at a single node
(one per number of laps) and a nontrivial endomorphism. The proofs of all
stated results have been machine-checked; we record only the mathematical
content.
\end{abstract}

\tableofcontents
\bigskip

%======================================================================
\section{Introduction}
%======================================================================

Directed containers of shipping networks sit on a ladder of structural
hypotheses. At the polytree rung, paths between nodes are unique and position
sets of finite graphs are finite. Climbing the ladder---first to directed
acyclic graphs that are not polytrees, then to graphs with directed
cycles---destroys these properties in turn.

The present note records two negative results that mark those rungs.
\begin{enumerate}
  \item \textbf{DAG rung.} On a directed acyclic graph obtained by adding a
        Singapore transshipment alternative to a direct Melbourne--Thailand sea
        leg, there are two distinct paths between the same endpoints. Hence
        $\mathrm{network}(G)$ is not thin.
  \item \textbf{Graph rung.} On a two-node directed cycle modelling a scrap
        loop between a rolling mill and a melt shop, the set of positions at
        the rolling mill is infinite: the map sending $n \in \N$ to the path
        that traverses the cycle $n$ times is injective. The one-lap path is a
        nontrivial endomorphism (distinct from the root).
\end{enumerate}

The development works throughout in a single universe of types. Proofs are
indicated by the phrase ``by lean,'' meaning they have been checked by a proof
assistant rather than written out in prose.

\paragraph{What this note claims.}
In ordinary mathematical language: acyclicity alone does not imply unique
paths, and the presence of a directed cycle yields infinitely many endomorphisms
at a node of the free category. Both phenomena are witnessed by fully explicit
finite graphs and fully explicit paths.

%======================================================================
\section{Preliminaries}
%======================================================================

We work with directed graphs and their path types as in the network
development: nodes~$V$, edges~$E$, endpoint maps $\src$ and $\tgt$, and
inductively defined finite directed paths with concatenation. The network
directed container $\mathrm{network}(G)$ has shapes $V$ and positions at~$v$
given by endpoint--path pairs $\sum_w \mathrm{Path}_G(v,w)$.

\begin{definition}[Path length]\label{def:path-len}
The \textbf{length} $\lvert p\rvert$ of a path~$p$ is the number of edges it
contains: $\lvert\mathrm{nil}(v)\rvert = 0$ and
$\lvert\mathrm{cons}(e,p)\rvert = \lvert p\rvert + 1$.
\end{definition}

Length is a convenient invariant for distinguishing paths without comparing
them edge by edge.

%======================================================================
\section{The DAG rung: two routes via Singapore}
%======================================================================

We extend the eleven-facility Sawasdee Cellars shipping graph by a Singapore
transshipment option between Melbourne Port and Thai Port.

\begin{definition}[Extended Sawasdee graph]\label{def:G-prime}
Let $G'$ be the directed graph whose nodes are the original eleven facilities
together with an additional node $\mathrm{singapore}$, and whose edges are the
original ten shipping legs together with two new legs
\begin{align*}
  \mathrm{sing}_1 &\;:\; \mathrm{melbournePort} \to \mathrm{singapore}, \\
  \mathrm{sing}_2 &\;:\; \mathrm{singapore} \to \mathrm{thaiPort}.
\end{align*}
The original direct sea leg
$\mathrm{sea} : \mathrm{melbournePort} \to \mathrm{thaiPort}$ is retained. The
graph $G'$ is acyclic (it remains a DAG) but is no longer a polytree: the two
outgoing edges $\mathrm{sea}$ and $\mathrm{sing}_1$ from Melbourne Port reconverge
at Thai Port.
\end{definition}

\begin{definition}[Two paths]\label{def:two-paths}
Define paths in $G'$ from $\mathrm{melbournePort}$ to $\mathrm{thaiPort}$ by
\begin{align}
  \rho_{\mathrm{direct}}
  &\;=\;
  \mathrm{cons}\bigl(\mathrm{sea},\, \mathrm{nil}(\mathrm{thaiPort})\bigr),
  \\[4pt]
  \rho_{\mathrm{transship}}
  &\;=\;
  \mathrm{cons}\bigl(\mathrm{sing}_1,\,
    \mathrm{cons}\bigl(\mathrm{sing}_2,\,
      \mathrm{nil}(\mathrm{thaiPort})\bigr)\bigr).
\end{align}
\end{definition}

\begin{proposition}[Lengths]\label{prop:dag-lengths}
One has $\lvert\rho_{\mathrm{direct}}\rvert = 1$ and
$\lvert\rho_{\mathrm{transship}}\rvert = 2$.
\end{proposition}

\begin{proof}
By lean.
\end{proof}

\begin{theorem}[DAG rung: failure of thinness]\label{thm:dag-rung}
The type $\mathrm{Path}_{G'}(\mathrm{melbournePort},\, \mathrm{thaiPort})$ is
not a subsingleton. In particular, $\mathrm{network}(G')$ is not thin at
$\mathrm{melbournePort}$.
\end{theorem}

\begin{proof}
By lean.
\end{proof}

\begin{remark}
If the two paths were equal, they would have equal length, contradicting
Proposition~\ref{prop:dag-lengths}. Thus distinct shipping alternatives between
the same ports---a commonplace in real logistics---immediately leave the thin
rung of the ladder.
\end{remark}

%======================================================================
\section{The graph rung: an infinite scrap loop}
%======================================================================

\begin{definition}[Scrap-loop digraph]\label{def:scrap}
Let $G_{\mathrm{scrap}}$ be the directed graph with two nodes
$\{\mathrm{rollingMill},\, \mathrm{meltShop}\}$ and two edges
\begin{align*}
  \mathrm{rm2ms} &\;:\; \mathrm{rollingMill} \to \mathrm{meltShop}, \\
  \mathrm{ms2rm} &\;:\; \mathrm{meltShop} \to \mathrm{rollingMill}.
\end{align*}
\end{definition}

\begin{proposition}[Length of a concatenation]\label{prop:len-append}
For paths $p : \mathrm{Path}_{G_{\mathrm{scrap}}}(v,w)$ and
$q : \mathrm{Path}_{G_{\mathrm{scrap}}}(w,u)$,
\[
  \lvert p \mathbin{{+}{+}} q\rvert \;=\; \lvert p\rvert + \lvert q\rvert.
\]
\end{proposition}

\begin{proof}
By lean.
\end{proof}

\begin{definition}[Laps]\label{def:laps}
Write
\begin{equation}
  \lambda
  \;=\;
  \mathrm{cons}\bigl(\mathrm{rm2ms},\,
    \mathrm{cons}\bigl(\mathrm{ms2rm},\,
      \mathrm{nil}(\mathrm{rollingMill})\bigr)\bigr)
\end{equation}
for one full lap of the cycle (length~$2$). Define
$\mathrm{laps} : \N \to \mathrm{Path}_{G_{\mathrm{scrap}}}
(\mathrm{rollingMill},\, \mathrm{rollingMill})$ by
\begin{equation}
  \mathrm{laps}(0) \;=\; \mathrm{nil}(\mathrm{rollingMill}),
  \qquad
  \mathrm{laps}(n+1) \;=\; \mathrm{laps}(n) \mathbin{{+}{+}} \lambda.
\end{equation}
\end{definition}

\begin{proposition}[Length of $n$ laps]\label{prop:laps-len}
For every $n \in \N$,
\[
  \lvert\mathrm{laps}(n)\rvert \;=\; 2n.
\]
\end{proposition}

\begin{proof}
By lean.
\end{proof}

\begin{proposition}[Injectivity of laps]\label{prop:laps-inj}
The map $\mathrm{laps}$ is injective: if $\mathrm{laps}(m) = \mathrm{laps}(n)$,
then $m = n$.
\end{proposition}

\begin{proof}
By lean.
\end{proof}

\begin{definition}[Infinite type]\label{def:infinite}
A type $\alpha$ is \textbf{infinite} when there exists an injection
$\N \hookrightarrow \alpha$.
\end{definition}

\begin{theorem}[Graph rung: infinite positions]\label{thm:graph-rung}
The position set of $\mathrm{network}(G_{\mathrm{scrap}})$ at
$\mathrm{rollingMill}$ is infinite. Concretely, the map
\begin{equation}
  n \;\longmapsto\;
  \bigl\langle \mathrm{rollingMill},\, \mathrm{laps}(n) \bigr\rangle
\end{equation}
is an injection
$\N \hookrightarrow
\sum_w \mathrm{Path}_{G_{\mathrm{scrap}}}(\mathrm{rollingMill},\, w)$.
\end{theorem}

\begin{proof}
By lean.
\end{proof}

\begin{definition}[Nontrivial endomorphism]\label{def:nontrivial-endo}
Let
\begin{equation}
  \eta
  \;=\;
  \bigl\langle \mathrm{rollingMill},\, \lambda \bigr\rangle
  \;\in\;
  \Pos\bigl(\mathrm{network}(G_{\mathrm{scrap}})\bigr)
    (\mathrm{rollingMill})
\end{equation}
be the one-lap position at the rolling mill.
\end{definition}

\begin{theorem}[Nontrivial endomorphism]\label{thm:nontrivial-endo}
The position $\eta$ is distinct from the root position at
$\mathrm{rollingMill}$.
\end{theorem}

\begin{proof}
By lean.
\end{proof}

\begin{remark}
Length distinguishes them: $\lvert\lambda\rvert = 2$ while the empty path has
length~$0$. In the free category on $G_{\mathrm{scrap}}$, $\eta$ is a
non-identity endomorphism of the rolling-mill object; iterating it produces the
infinite family of Theorem~\ref{thm:graph-rung}.
\end{remark}

%======================================================================
\section{Discussion}
%======================================================================

\paragraph{What has been established.}
Two concrete witnesses separate the rungs of the structural ladder for network
directed containers.
\begin{itemize}
  \item On a DAG that is not a polytree, multiple paths between the same
        endpoints exist, so thinness fails.
  \item On a graph with a directed cycle, the position set at a node on the
        cycle is infinite, and nontrivial endomorphisms exist.
\end{itemize}
Both conclusions are elementary path-counting facts; their role is to mark
sharp boundaries for the positive theorems available on the polytree rung.

\paragraph{Scope.}
The development does not classify all DAGs or all cyclic graphs, nor does it
quantify the growth of position sets beyond the injection from~$\N$. The
examples are minimal by design.

\paragraph{Machine checking.}
Every definition and theorem statement above corresponds to a declaration in a
machine-checked formalisation. The phrase ``by lean'' indicates that the proof
has been verified by a proof assistant.

%======================================================================
\section*{Acknowledgements}
%======================================================================

The mathematical content follows the DAG and graph rungs of the supply-chain
development that motivated this note. The presentation is a prose rendering of
a formal development intended for human review of the theorem statements.

\begin{thebibliography}{9}

\bibitem{abbott2005}
M.~Abbott, T.~Altenkirch, and N.~Ghani.
\newblock Containers: constructing strictly positive types.
\newblock \emph{Theoretical Computer Science}, 342(1):3--27, 2005.

\bibitem{ahman2014}
D.~Ahman and T.~Uustalu.
\newblock Coalgebraic update lenses.
\newblock In \emph{MFPS}, volume~308 of \emph{ENTCS}, pages 25--48, 2014.

\end{thebibliography}

\end{document}
